% GENERATED by CANON/book_toolbox.py from CANON/catalogue.json and CANON/instruments.json. Do not edit by hand.
\chapter{The toolbox}\label{ch:toolbox}

Every reading in this part comes from one of a small set of tools. Each answers its own question and has its own specialty; together they look at a sample from different sides, so that what one cannot see another can. This chapter introduces each tool and says where it stands: \emph{commissioned} (it passed every test written for it before its data were read, and its readings in that scope are results), \emph{in development}, or a \emph{candidate} still being designed. The chapters that follow give the physics and the evidence.

\section{The instruments}

\subsection*{IAM-Atlas v2.1}
\textbf{The question.} Whose DNA is this? Which cells are in a sample, and in what shares?\par
\textbf{How it reads.} A map of 79 kinds of healthy human cells, built from seven public sources on arrays and sequencing, fitted as one model so every spot on the genome carries its value and its uncertainty. Like a state estimator across the whole interconnection: many utilities' meters, each weighted by how far it can be trusted.\par
\textbf{Its specialty.} It is the map every other instrument reads against: for 79 healthy cell types, the expected tagging at every site with its range of error, so a sample can be split into its cells and each cell compared with its own healthy kind. (Chapter~\ref{ch:atlas})\par
\textbf{Reads from.} arrays or sequencing. \textbf{Status.} Commissioned (three cell references (450K) and the colon and lung blind sites, 2026-10-16).\par

\subsection*{Met-A}
\textbf{The question.} How far has a tissue's tagging level drifted toward the ceiling?\par
\textbf{How it reads.} Reads the methylation level at sites a healthy cell holds firmly, against a healthy reference, on cheap and plentiful arrays. Commissioned on isolated neutrophils; whole blood is in development.\par
\textbf{Its specialty.} Cost and reach: a single array, the most common methylation measurement in the world, read on the gauge against the same cell type's healthy reference, with its detection limit printed. (Chapter~\ref{ch:meta})\par
\textbf{Reads from.} arrays. \textbf{Status.} Commissioned (isolated neutrophils, EPIC v1, 2026-10-09).\par

\subsection*{C-Score (the clustering meter)}
\textbf{The question.} Has the pattern started to clump?\par
\textbf{How it reads.} Cancer does not only shift tagging levels; it clumps them into runs. The C-Score measures that clumping against a healthy band. The band held in three laboratories that never helped set it, and one person measured twice reads the same.\par
\textbf{Its specialty.} The shape of a change, not its size: whether a sample's departures from healthy are scattered like ordinary wear or packed together in stretches, the way a tumor that silences or loses a region leaves them. It is the clustering step of the cosmic microwave background analysis chain, carried over to the genome. (Chapter~\ref{ch:cscore})\par
\textbf{Reads from.} arrays (whole-genome sequencing too). \textbf{Status.} Commissioned (whole blood, EPIC v1, 2026-10-15).\par

\subsection*{IAM-A}
\textbf{The question.} Is each mark still held at its own healthy rate?\par
\textbf{How it reads.} Counts tagging mistakes one DNA molecule at a time, the most direct reading of informational fidelity we have, and reads them against the sample's own line. The physics end of the catalogue: the floor and ceiling are derived, and since 2026-10-15 the healthy middle comes from the sample itself.\par
\textbf{Its specialty.} That it reads single molecules and carries its own reference: how well each molecule held its marks, compared with the healthy level set by the cell's own loose share, with no library of other people's cells. (Chapter~\ref{ch:iama})\par
\textbf{Reads from.} single-molecule sequencing (bisulfite, EM-seq; nanopore next). \textbf{Status.} In development.\par

\subsection*{Group-A}
\textbf{The question.} Has a whole group moved (a trial arm, an exposure, a cohort)?\par
\textbf{How it reads.} Single blood readings are too noisy for one person on arrays, but a group average is not. Our honesty test showed two healthy groups from two labs can differ as much as a tumor would, so Group-A compares groups only when they were run in the same series.\par
\textbf{Its specialty.} Groups: whether a whole cohort, trial arm or exposure has moved, read as one, with patients and healthy comparisons run in the same laboratory series.\par
\textbf{Reads from.} arrays. \textbf{Status.} In development.\par

\section{The tools the instruments run on}

\begin{itemize}
\item \textbf{Intake and calibration.} Takes an array's raw intensities, checks every control, and turns them into tagging levels on one fixed scale, refusing what it cannot read and saying why. (Chapter~\ref{ch:instrument})
\item \textbf{The tare.} Reads each specimen against healthy references run beside it, so a laboratory's own offset cancels instead of passing for disease. (Chapter~\ref{ch:instrument})
\item \textbf{Composition.} For whole blood, works out which cells are in the sample and in what shares, so a change in the mix is not mistaken for a change in the cells. (Chapter~\ref{ch:chain})
\item \textbf{The deconvolution tool.} Reads a mixed sample against atlas v2.1 and reports each cell's share, after a sample check that refuses arrays it cannot read. (Chapter~\ref{ch:atlas})
\item \textbf{The sky map.} Lays the genome on a sphere, chromosome by chromosome, and draws each site's departure from healthy, the way a microwave sky is drawn; the same residuals the clustering meter scores. (Chapter~\ref{ch:skytools})
\item \textbf{The chain and its report.} Runs every stage in a fixed order, records every input and setting, and prints one report per sample with each instrument's reading, its status and its detection limit. (Chapter~\ref{ch:report})
\end{itemize}

\section{Candidates}

Readings still being designed, each with what it would add and what it cannot see. A candidate becomes an instrument only by passing tests written before its data are read.

\begin{itemize}
\item \textbf{Territory contrast.} \emph{What it reads:} where the loss happens first - a sustained energy shortfall should take the late-replicating, solo-WCGW territory before it touches housekeeping promoters \emph{What it adds:} self-taring - a ratio inside one array, so plate, batch and much of the composition effect cancel because numerator and denominator see the same cell mixture. Attacks the exact spread problem that blocked the blood answer tonight \emph{What it cannot see:} says nothing about absolute fidelity; a loss uniform across territories is invisible (confirmed, 0.1 \% of the signal). And the practical failure: a difference cancels an additive plate offset exactly but cannot cancel a multiplicative one - a 10 \% composition change moves it about 5x what a 2 \% shortfall does - while the beta = 1 ceiling makes offsets and noise one-sided
\item \textbf{Met-A on whole blood.} \emph{What it reads:} total disorder above what the estimated cell mixture explains - the existing composition-aware reading, in a tube of blood \emph{What it adds:} group comparisons are within reach now: with 50 people a mean loss under 1 \% separates the group from healthy, and the function already exists in the chain \emph{What it cannot see:} per-person detection in blood: the spread between healthy people (6.4 \%) is larger than the tissue-sized signal (2-7 \%), so one person cannot be called. An array also averages molecules before anyone reads them, so one cell type can never be isolated. And these arrays are spent as the C-score test, so this was development use And across laboratories without a shared series the raw reading carries a 3.7 \% offset, which is inside the tumor range
\item \textbf{Distance to the ceiling.} \emph{What it reads:} how close a sample sits to the derived entropy ceiling - the right-hand mark of the gauge, past which the reading stops rising with further loss \emph{What it adds:} derived with no free parameter, and it answers the author's own question - how far from the ceiling a tumor arrives \emph{What it cannot see:} only informative for samples near the ceiling; says nothing about a healthy-range sample
\item \textbf{Run-length spectrum of losses.} \emph{What it reads:} the shape of clustering - exponential against heavy-tailed, which separates a processive mechanism from independent point losses \emph{What it adds:} same data as the clustering meter, different question; distinguishes mechanism rather than only size \emph{What it cannot see:} impossible on arrays; needs depth, so shallow libraries are out
\item \textbf{Tail of fragments in plasma (the early-detection instrument).} \emph{What it reads:} the few badly disordered molecules in a plasma sample, not the average - because early cancer is a tiny minority of fragments and an average hides it by construction \emph{What it adds:} the highest reward on the whole board - it is the blood test, and a tail statistic is the right shape for early detection \emph{What it cannot see:} needs depth and fragment length; a shallow or heavily fragmented sample has no tail to read. Composition of plasma is itself a question
\end{itemize}

\section{Why several}

Each instrument answers a different question about the same sample, so together they leave no blind side. The atlas says whose DNA it is. Met-A says how far a tissue's tagging level has drifted toward the ceiling. The C-Score says whether the pattern has started to clump. IAM-A says whether each mark is still held at its own healthy rate. Group-A says whether a whole group has moved, for studies rather than single patients. A tumor that hides from one reading still has to answer to the others.

